\documentclass[11pt,letterpaper]{article}
\usepackage[T1]{fontenc}
\usepackage[utf8]{inputenc}
\usepackage{lmodern}
\usepackage[margin=1in]{geometry}
\usepackage{amsmath,amssymb,amsthm,mathtools}
\usepackage{booktabs,array,longtable,enumitem}
\usepackage{microtype,xcolor,graphicx}
\usepackage[hyphens]{url}
\usepackage[colorlinks=true,linkcolor=blue!45!black,citecolor=blue!45!black,urlcolor=blue!45!black]{hyperref}
\usepackage[nameinlink,noabbrev]{cleveref}
\hypersetup{pdftitle={Unsmoothed prime fluctuations and information loss in symplectic projections},pdfauthor={ChatGPT; prepared for Vladimir Reshetnikov},pdfsubject={Two extensions of the ProveIt GRH development},pdfkeywords={Dirichlet GRH, prime counting, symplectic matrices, sample complexity}}
\usepackage{fancyhdr}
\pagestyle{fancy}
\fancyhf{}
\fancyhead[L]{\small Two extensions of the ProveIt GRH development}
\fancyhead[R]{\small 27 November 2026}
\fancyfoot[C]{\thepage}
\setlength{\headheight}{14pt}
\setlength{\emergencystretch}{2em}
\setlength{\parskip}{3pt}
\setlist{itemsep=3pt,topsep=4pt}
\numberwithin{equation}{section}
\newtheorem{theorem}{Theorem}[section]
\newtheorem{lemma}[theorem]{Lemma}
\newtheorem{proposition}[theorem]{Proposition}
\newtheorem{corollary}[theorem]{Corollary}
\theoremstyle{definition}
\newtheorem{definition}[theorem]{Definition}
\theoremstyle{remark}
\newtheorem{remark}[theorem]{Remark}
\newcommand{\F}{\mathbb F}
\newcommand{\R}{\mathbb R}
\newcommand{\C}{\mathbb C}
\newcommand{\Z}{\mathbb Z}
\newcommand{\N}{\mathbb N}
\newcommand{\TV}{\operatorname{TV}}
\DeclareMathOperator{\rank}{rank}
\DeclareMathOperator{\Var}{Var}
\DeclareMathOperator{\Cov}{Cov}
\DeclareMathOperator{\Span}{span}
\DeclareMathOperator{\diag}{diag}
\newcommand{\E}{\mathbb E}
\newcommand{\Prob}{\mathbb P}
\newcommand{\ind}{\mathbf 1}
\title{\textbf{Unsmoothed prime fluctuations and information loss in symplectic projections}\\[8pt]
\large Two proved extensions of the ProveIt GRH development}
\author{Research note prepared for Vladimir Reshetnikov\\\small Mathematical development and source review by ChatGPT}
\date{27 November 2026}
\begin{document}
\maketitle
\begin{abstract}
We develop two extensions of the supplied ProveIt GRH sources, pinned to commit
\texttt{3d8689ceae76}.
First, using the source's Dirichlet GRH statement, we prove strict positive
definiteness of the covariance of ordinary, unsmoothed prime-counting errors
in all reduced residue classes at any finite set of distinct multiplicative
scales. The result survives every finite deletion of zero frequencies and
gives positive occupation of both directional tails in every sufficiently
long logarithmic interval. It uses no simplicity or linear independence of
zero ordinates and settles the unsmoothed multiscale clause of the article's
explicit question on less smoothing. The analytic bridge is convergence in a
uniform local fourth-moment norm, together with a controlled transfer from
prime powers to actual primes.

Second, we prove an unconditional statistical separation for rectangular
symplectic projections over every finite field. For an $n\times(n+c)$
projection with nonsingular source form, $1\le c\le n$ and $n+c$ even,
the full matrix laws become perfectly distinguishable as $c\to\infty$,
whereas their rank laws have distance asymptotic to $q^{-n}/(q-1)$.
We compute their rank chi-square divergence and prove that distinguishing
from ranks alone requires order $q^{2n-c}$ independent observations.
The proofs include exact finite formulas and cover the endpoint $c=n$.
An exact-arithmetic verification program and a source audit accompany the paper.
These are new extensions relative to the inspected article; external priority
for their complete formulations is not certified by the limited literature search.
\end{abstract}
\tableofcontents
\newpage

\section{Source analysis, dependencies, and contributions}\label{sec:context}

\subsection{What was inspected}
The two requested directories were read at the fixed ProveIt commit
\texttt{3d8689ceae76}.
The vendor directory identifies its upstream snapshot as
\texttt{lsernau/grh-lean}, commit
\texttt{72fead102f66}.
The archive records both full commit identifiers. Its Git tree contains 818 files. The article README describes a merge of
120 reports, with 689 pages and 869 numbered results; it explicitly says the
merge is not an independent mathematical review \cite{ProveItArticle}.
We materialized and indexed the article's 21 section files, inspected the
relevant arguments in its source, zero-fluctuation, and spanning sections,
and inspected the vendored theorem interfaces and the principal reduction
and spanning files. A complete list of the source paths and expected Git blob
hashes used in this review is included in the archive.

The exact public Lean interface uses Mathlib's Riemann zeta and Dirichlet
$L$-functions. The file \texttt{src/Statements.lean} imports only Mathlib;
\texttt{src/Main.lean} exports \texttt{riemannHypothesis},
\texttt{riemannZeta\_eq\_zero\_re},
\texttt{LFunction\_ne\_zero\_of\_half\_lt\_re},
\texttt{LFunction\_eq\_zero\_re}, and the primitive-character specialization
\cite{ProveItVendor}. For our arithmetic theorem the input is precisely this:
every nontrivial zero of a primitive Dirichlet $L$-function lies on
$\Re s=1/2$. This includes the Riemann zeta function. It asserts neither
simple zeros nor nonvanishing at the central point. It is not an interface
for general nonabelian Artin or automorphic GRH.

The package README reports an axiom check and independent proof replay.
Those reports were read, not rerun here: the workspace has no Lean or Lake
executable. The vendoring README also says the separate package is not part
of ProveIt's root Lake build. The five \texttt{sorry} placeholders in
\texttt{Challenge.lean} are declared comparator targets and are not imports
of the inspected proof path. Their presence alone is not evidence of a hole
in the exported theorem. Our audit is a scoped source review, not a fresh
kernel verification. The new proofs below are human-readable proofs and have
not been formalized in Lean.

\subsection{The mathematical mechanisms}
The source's arithmetic argument compares the same covolume and arithmetic
index from above and below. A near-complete spanning family permits
cancellation of the large main terms. In the notation of the inspected
\texttt{star\_inequality} declaration, the remainder satisfies a bound of
the form $mR\le f(a)(\log b)^5$. After normalization this becomes an eventual
one-sided prime-remainder estimate for every exponent above $1/2$.
An auxiliary quadratic twist removes unwanted zeros in a prescribed finite
rectangle. A Landau-type argument applied to the remainder's transform then
forces the remaining nonnegative zero multiplicities to vanish in the
half-plane to the right of $1/2$.

Two parts of that architecture motivate the present paper. Its zero-detection
method makes arithmetic prime probes useful for recovering coefficients of
spectral combinations. The article turns these probes into a multiscale
spanning theorem for zero-multiplicity vectors. Its finite spanning method
also uses additive characters, positive-type weights, and singularity
estimates for compressions of group actions. The article develops this into
exact symplectic compression laws. We carry the first route to unsmoothed
prime counts and the second to a sharp loss of statistical information.

\subsection{What is added}
The relevant article statements are identified by their stable LaTeX labels,
so they can be located despite changes in pagination.
\begin{center}\small
\begin{tabular}{>{\raggedright\arraybackslash}p{0.28\textwidth}>{\raggedright\arraybackslash}p{0.29\textwidth}>{\raggedright\arraybackslash}p{0.34\textwidth}}
\toprule
Existing result & Extension proved here & Dependency\\
\midrule
\texttt{zer:thm:coercive}, \texttt{zer:cor:multiscale}: smoothed multiscale covariance
& Ordinary prime counts, correct prime-square and central-zero centering,
finite spectral deletions, uniform directional occupation
& Dirichlet GRH; standard explicit formulas\\[5pt]
\texttt{spn:thm:complete-law}, \texttt{spn:thm:rank-blind}: exact compression law and square rank obstruction
& Wide rectangular limit, boundary constants, uniform chi-square asymptotic,
sharp order of rank-only sample complexity
& Unconditional finite-field linear algebra\\
\bottomrule
\end{tabular}
\end{center}

The first result answers the multiscale unsmoothed-count clause of item 10
in \texttt{zer:sec:open}; it does not assert all possible orderings of prime
races or resolve every question about unsmoothed derivatives.
The second gives a concrete growing-dimensional family for
\texttt{spn:q:symplectic}. The exact rectangular likelihood is a specialization
and simplification of the article's general count. The additional mathematical
content is the uniform asymptotic analysis and the testing threshold, not a
claim to have discovered the prescribed-Gram orbit count.

Stepanov almost periodicity of normalized prime errors is classical
\cite{KaczorowskiRamare2003}; general limiting-distribution methods also precede
this work \cite{AkbaryNgShahabi2014}. Results asserting all orderings of prime
races are stronger and may use extra independence assumptions
\cite{MartinNg2019}. Finite-field rank laws and prescribed-form sampling likewise
have established theories \cite{FulmanGoldstein2015,Ishizuka2026}. The present
novelty claim is therefore scoped to the two precise extensions just stated,
relative to the supplied article. The accompanying novelty record describes
the targeted primary-source comparison; it does not establish worldwide priority.

\section{Unsmoothed prime counts at arbitrary multiplicative scales}
\label{sec:pf}

The supplied article proves positive definite covariance for smoothed prime
errors at arbitrary finitely many multiplicative scales.  Its open questions
ask whether this persists for unsmoothed prime counts after the prime-square
bias is removed \cite[Section~\textup{``Zeros, Euler products and prime races''},
open question~10]{ProveItArticle}.  We prove the multiscale part of that
extension for the ordinary counting function $\pi(x;q,a)$.  The conclusion
survives removal of every fixed finite set of zero frequencies, and yields
positive and negative excursions of every centered linear contrast in
\emph{every sufficiently long logarithmic interval}.

The only hypothesis about zero locations is Dirichlet GRH, used in the precise
sense supplied by the vendored development \cite{ProveItVendor}.  Simplicity,
disjointness of different zero sets, and rational independence of ordinates
are not assumed.  Stepanov almost periodicity of normalized prime-counting
remainders is classical: Kaczorowski and Ramar\'e prove an $L^2$ version in a
more general setting \cite[Theorem~1(e)]{KaczorowskiRamare2003}.  The extension
here combines an explicit local $L^4$ argument, a transfer to actual prime
counts, and the arithmetic separation of arbitrary scales.  We do not claim
that every ordering of the uncentered prime counts occurs, or settle the
separate question about all jet derivatives.

\subsection{Statement and normalizations}
\label{subsec:pf-statement}

Fix $q\geq1$, write $G=(\mathbb Z/q\mathbb Z)^\times$ and
$n=\varphi(q)$, and use the trivial group and $n=1$ when $q=1$.
For a character $\chi$ modulo $q$, let $\chi^*$ be its primitive inducer and
let $m_\chi(\gamma)$ denote the multiplicity of the zero
$\tfrac12+i\gamma$ of $L(s,\chi^*)$.  A missing zero has multiplicity zero.
Let $\Gamma_q$ be the set of distinct ordinates occurring in this finite
family, and define
\begin{equation}
\label{eq:pf-v}
 v_a(\gamma)=\frac1n\sum_{\chi\bmod q}
       \overline{\chi(a)}m_\chi(\gamma),
 \qquad
 r_q(a)=\#\{b\in G:b^2=a\}.
\end{equation}
Thus coincidences between zeros of different functions are grouped before
any square is taken.  Conjugation gives
$v_a(-\gamma)=\overline{v_a(\gamma)}$, and $v_a(0)$ is real.
Put $\operatorname{li}(x)=\int_2^x du/\log u$, and set
\begin{equation}
\label{eq:pf-X}
 X_a(t)=t e^{-t/2}
   \left(\pi(e^t;q,a)-\frac{\operatorname{li}(e^t)}n\right)
       +\frac{r_q(a)}n+2v_a(0),\qquad t\geq\log2.
\end{equation}
The term $r_q(a)/n$ removes prime squares, while $2v_a(0)$ removes the
possible central zero contribution.  The latter is not assumed to vanish.
These constants differ from those of the article's smoothed ramp statistic
because the Mellin denominator is now $s$, rather than $s(s-1)$.

Fix distinct real shifts $h_1,\ldots,h_r$ and put $D=rn$.
They correspond to the multiplicative scales $e^{h_1},\ldots,e^{h_r}$.
For $\gamma\ne0$ define a vector in $\mathbb C^D$ by
\begin{equation}
\label{eq:pf-z}
 z(\gamma)_{j,a}
       =-\frac{e^{i\gamma h_j}v_a(\gamma)}{\tfrac12+i\gamma}.
\end{equation}
If $J\subset\Gamma_q\setminus\{0\}$ is a finite symmetric set, remove its
frequencies from the observed vector:
\begin{equation}
\label{eq:pf-deletion}
 \mathcal X_J(t)_{j,a}
       =X_a(t+h_j)-\sum_{\gamma\in J}
                         z(\gamma)_{j,a}e^{i\gamma t}.
\end{equation}
The vector is real.  Its domain starts at any fixed $u_0$ for which
$u_0+h_j\geq\log2$ for every $j$.

\begin{theorem}[Unsmoothed multiscale covariance and occupation]
\label{thm:pf-main}
Assume Dirichlet GRH for the characters modulo $q$.  For the data just
defined, the series
\begin{equation}
\label{eq:pf-cov}
 C_J=\sum_{\gamma\in\Gamma_q\setminus(J\cup\{0\})}
                   z(\gamma)z(\gamma)^*
\end{equation}
converges to a real symmetric positive definite matrix.  Uniformly over
$u\geq u_0$, as $L\longrightarrow\infty$,
\begin{equation}
\label{eq:pf-moments}
 \frac1L\int_u^{u+L}\mathcal X_J(t)\,dt\longrightarrow0,
 \qquad
 \frac1L\int_u^{u+L}\mathcal X_J(t)\mathcal X_J(t)^{\mathsf T}\,dt
                       \longrightarrow C_J.
\end{equation}
There are constants $\varepsilon,\eta,L_0>0$ such that, for every real unit
vector $w\in\mathbb R^D$, every $u\geq u_0$ and every $L\geq L_0$,
\begin{align}
\label{eq:pf-positive}
 \bigl|\{t\in[u,u+L]:w\cdot\mathcal X_J(t)\geq\varepsilon\}\bigr|
       &\geq\eta L,\\
\label{eq:pf-negative}
 \bigl|\{t\in[u,u+L]:w\cdot\mathcal X_J(t)\leq-\varepsilon\}\bigr|
       &\geq\eta L.
\end{align}
Here $|\cdot|$ denotes Lebesgue measure.  Consequently the closed convex
hull of the essential values on each such interval contains the closed
Euclidean ball of radius $\varepsilon$ about zero.

Moreover, there are $c,K,T_0>0$, depending only on $q$ and the shifts,
such that
\begin{equation}
\label{eq:pf-band}
 \lambda_{\min}(C_J)\geq\frac{cH}{T^2\log(2T)}
\end{equation}
whenever $T\geq T_0$, $K\log(2T)\leq H\leq T$, and
$T>\max_{\gamma\in J}|\gamma|$; the last maximum is zero if $J$ is empty.
\end{theorem}

The constants may depend on the fixed shifts and deteriorate when shifts
coalesce.  No uniformity as the number of coordinates grows is asserted.
In terms of $x=e^t$, the intervals above are $[e^u,e^{u+L}]$ and their
measure is $dx/x$.  Thus the theorem supplies bounded waiting on the
logarithmic scale for every centered directional contrast, with a positive
proportion of logarithmic measure on both sides.

\subsection{Local fourth moments of zero sums}
\label{subsec:pf-local}

For a locally $p$-integrable function on $\mathbb R$, write
\[
 \|f\|_{S^p}=\sup_{u\in\mathbb R}
           \left(\int_u^{u+1}|f(t)|^p\,dt\right)^{1/p}.
\]
For vector-valued functions, use the Euclidean norm inside the integral.
Convergence in this norm is uniform local $L^p$ convergence; in particular
it implies convergence as distributions.  We first prove the analytic
estimate that permits unsmoothing.

\begin{lemma}[Binned local moments]
\label{lem:pf-bins}
Let $P(t)=\sum_\lambda c_\lambda e^{i\lambda t}$ be a finite exponential
sum with real frequencies, and put
$b_j=\sum_{\lambda\in[j,j+1)}|c_\lambda|$.  There is an absolute constant
$A$ such that
\begin{equation}
\label{eq:pf-bins}
 \|P\|_{S^2}\leq A\|b\|_{\ell^2},
 \qquad
 \|P\|_{S^4}\leq A\|b\|_{\ell^{4/3}}.
\end{equation}
Consequently, if a countable locally finite frequency family has
$b\in\ell^{4/3}(\mathbb Z)$, its finite symmetric truncations converge in
$S^4$.
\end{lemma}

\begin{proof}
Consider the nonnegative function
\[
 W(s)=\left(\frac{\sin(s/2)}{s/2}\right)^2,
 \qquad W(0)=1.
\]
It has a positive lower bound on $[-1/2,1/2]$ and satisfies the exact
identity
\[
 W(s)=\int_{-1}^1(1-|\xi|)e^{i\xi s}\,d\xi.
\]
Its Fourier transform is therefore bounded and supported in $[-1,1]$,
with the usual factor $2\pi$ in the inverse convention.  Multiplying
$|P|^2$ by $W(t-u-1/2)$ and integrating over $\mathbb R$ gives
\begin{align*}
 \int_u^{u+1}|P(t)|^2\,dt
 &\leq C\int_{\mathbb R}|P(t)|^2W(t-u-1/2)\,dt\\
 &\leq C'\sum_{|\lambda-\mu|\leq1}
                             |c_\lambda||c_\mu|\\
 &\leq C''\sum_j b_j^2.
\end{align*}
Only equal or adjacent bins can interact, and
$2b_jb_{j+1}\leq b_j^2+b_{j+1}^2$.  The constants do not depend on $u$.

Apply this $S^2$ estimate to $P^2$.  Its absolute coefficient mass $d_k$
in bin $[k,k+1)$ is at most
\[
 d_k\leq(b*b)_k+(b*b)_{k-1}.
\]
Indeed, adding frequencies in bins $j$ and $\ell$ gives a frequency in
bin $j+\ell$ or $j+\ell+1$; grouping coincident sums can only decrease
absolute mass.  Young's convolution inequality yields
\[
 \|P\|_{S^4}^4
   =\|P^2\|_{S^2}^2
   \leq C\|d\|_{\ell^2}^2
   \leq4C\|b*b\|_{\ell^2}^2
   \leq4C\|b\|_{\ell^{4/3}}^4.
\]
This proves the second estimate after enlarging $A$.  Applying it to
differences of truncations proves the final assertion, since the
$\ell^{4/3}$ norm of their omitted bin masses tends to zero.
\end{proof}

For the fixed Dirichlet family, the standard local zero count gives
\begin{equation}
\label{eq:pf-localcount}
 \sum_{\gamma\in\Gamma_q\cap[j,j+1)}\sum_{\chi\bmod q}m_\chi(\gamma)
                 \ll_q\log(2+|j|).
\end{equation}
This count includes multiplicities and holds for both signs.  Therefore
the absolute coefficient masses of
\begin{equation}
\label{eq:pf-F}
 F_a(t)=-\sum_{\substack{\gamma\in\Gamma_q\\\gamma\ne0}}
                  \frac{v_a(\gamma)}{\tfrac12+i\gamma}e^{i\gamma t}
\end{equation}
obey
\begin{equation}
\label{eq:pf-bdecay}
 b_j\ll_q\frac{\log(2+|j|)}{1+|j|}.
\end{equation}
The finitely many bins near zero cause no difficulty because
$|\tfrac12+i\gamma|\geq1/2$.
Lemma~\ref{lem:pf-bins} proves convergence of \eqref{eq:pf-F} in $S^4$,
and in particular a uniform bound for its local fourth moments.  Comparing
the tails of $\sum b_j^{4/3}$ and $\sum b_j^2$ with integrals also gives,
for $Y\geq2$,
\begin{equation}
\label{eq:pf-truncation}
 \left\|F_a+
   \sum_{0<|\gamma|\leq Y}
       \frac{v_a(\gamma)e^{i\gamma t}}{\tfrac12+i\gamma}\right\|_{S^4}
       \ll_q Y^{-1/4}\log(2Y),
\end{equation}
and the corresponding $S^2$ bound is
$O_q(Y^{-1/2}\log(2Y))$.  No lower bound on gaps between ordinates has
entered these estimates.

\subsection{From the explicit formula to ordinary prime counts}
\label{subsec:pf-transfer}

We use the classical explicit formula and local zero count for primitive
Dirichlet $L$-functions, in the distributional normalization also used in
\cite{KaczorowskiRamare2003}; the fixed-argument Landau formula needed
below is recorded separately.  The standard GRH consequence
\begin{equation}
\label{eq:pf-pointwise}
 \theta(x;q,a)=\frac{x}{n}+O_q\bigl(\sqrt{x}\log^2(2x)\bigr)
\end{equation}
will be used only for deterministic remainder estimates.  These are
classical analytic inputs, not new conclusions about zero locations.

Let
\[
 \psi(x;q,a)=\sum_{\substack{p^k\leq x\\p^k\equiv a\bmod q}}\log p,
 \qquad
 \theta(x;q,a)=\sum_{\substack{p\leq x\\p\equiv a\bmod q}}\log p.
\]

\begin{lemma}[The unsmoothed expansion]
\label{lem:pf-transfer}
For each $a\in G$, there is a remainder $R_a$ such that, almost everywhere
on $t\geq\log2$,
\begin{equation}
\label{eq:pf-XF}
 X_a(t)=F_a(t)+R_a(t),
 \qquad
 \|R_a\|_{L^4([U,U+1])}\ll_q U^{-1}\quad(U\geq2).
\end{equation}
The remainder is locally bounded, belongs to $L^2([\log2,\infty))$, and
has uniformly bounded local fourth moments on that half-line.
\end{lemma}

\begin{proof}
For a primitive character $\chi$, put $\delta_\chi=1$ in the conductor-one
case and $0$ otherwise.  The explicit formula, interpreted as an identity
of distributions on $t>0$, states
\begin{equation}
\label{eq:pf-psi-character}
 e^{-t/2}\bigl(\psi(e^t,\chi)-\delta_\chi e^t\bigr)
 =-\sum_\gamma\frac{m_\chi(\gamma)e^{i\gamma t}}
                         {\tfrac12+i\gamma}
       +O_\chi\bigl((1+t)e^{-t/2}\bigr),\qquad t\geq\log2.
\end{equation}
The term denoted by $O_\chi$ consists of the constant, the contribution
at zero, and the trivial zeros.  It is smooth away from $t=0$ and has
the stated bound.  One may obtain the distributional formula by
integrating $-L'/L(s,\chi)$ against a smooth compactly supported Mellin
test and moving the contour; rapid decay of the test transform makes
the zero sum convergent.  Dirichlet GRH turns each nontrivial zero term
into the pure phase displayed in \eqref{eq:pf-psi-character}.

The bin estimates prove $S^4$, hence $L^1_{\mathrm{loc}}$, convergence
of this same zero sum.  Its limit represents the distribution in
\eqref{eq:pf-psi-character}, so the equality holds almost everywhere
for the usual arithmetic step function.  The convention at its
countably many jumps does not affect any conclusion here.  If $\chi$
is imprimitive, the finitely many removed prime-power sequences have
total size $O_q(1+t)$ before normalization.  Thus the same estimate
holds with the zeros of $\chi^*$.

Character orthogonality and the central term
$-(\tfrac12)^{-1}v_a(0)=-2v_a(0)$ now give
\begin{equation}
\label{eq:pf-psi}
 E_{\psi,a}(t):=e^{-t/2}\left(\psi(e^t;q,a)-\frac{e^t}{n}\right)
       =F_a(t)-2v_a(0)+O_q\bigl((1+t)e^{-t/2}\bigr).
\end{equation}
In particular $E_{\psi,a}$ has bounded local fourth moments uniformly
on $t\geq\log2$.

The contribution of prime squares is
\[
 \psi(x;q,a)-\theta(x;q,a)
       =\sum_{b^2=a}\theta(\sqrt{x};q,b)
                  +O_q\bigl(x^{1/3}\log^2(2x)\bigr).
\]
The error follows by bounding the contributions of prime powers of
exponent at least three.  Applying \eqref{eq:pf-pointwise} at
$\sqrt{x}$, and then dividing by $\sqrt{x}$, yields
\begin{equation}
\label{eq:pf-theta}
 E_{\theta,a}(t):=e^{-t/2}\left(\theta(e^t;q,a)-\frac{e^t}{n}\right)
 =E_{\psi,a}(t)-\frac{r_q(a)}n
                    +O_q\bigl((1+t)^2e^{-t/6}\bigr).
\end{equation}
This also has uniformly bounded local fourth moments.  Notice that
the correction is $r_q(a)/n$, not $(r_q(a)-1)/n$: we are centering each
coordinate against $\operatorname{li}(x)/n$, including the principal
direction.

Partial summation gives the exact identity
\begin{align}
\label{eq:pf-volterra}
 t e^{-t/2}\left(\pi(e^t;q,a)-\frac{\operatorname{li}(e^t)}n\right)
 &=E_{\theta,a}(t)+V_a(t)+\frac{2t e^{-t/2}}{n\log2},\\
 V_a(t)&=t e^{-t/2}\int_{\log2}^t
                    \frac{e^{v/2}E_{\theta,a}(v)}{v^2}\,dv.
 \nonumber
\end{align}
We claim that $\|V_a\|_{L^4([U,U+1])}\ll_q U^{-1}$.
Split the integral at $v=t/2$.  On the earlier part,
\eqref{eq:pf-pointwise} gives $E_{\theta,a}(v)\ll_q(1+v)^2$, and the
resulting contribution is $O_q(t^2e^{-t/4})$.  On the later part put
$s=t-v$ and use $t/(t-s)^2\leq4/t$.  For $t\in[U,U+1]$,
\[
 |V_{a,\mathrm{late}}(t)|
 \leq\frac4U\int_0^{(U+1)/2}e^{-s/2}|E_{\theta,a}(t-s)|\,ds.
\]
For large $U$, all translated unit intervals on the right lie in the
half-line where the local fourth moments are bounded.  Minkowski's
inequality, followed by $\int_0^\infty e^{-s/2}ds=2$, proves the claim.

Combining \eqref{eq:pf-psi}, \eqref{eq:pf-theta}, and
\eqref{eq:pf-volterra} proves \eqref{eq:pf-XF} with exactly the
centering in \eqref{eq:pf-X}.  The remainders constructed in these
steps are locally bounded.  Finally
\[
 \int_U^{U+1}|R_a(t)|^2\,dt
       \leq\left(\int_U^{U+1}|R_a(t)|^4\,dt\right)^{1/2}
       \ll_q U^{-2},
\]
and summation over integer $U$ proves global square integrability.
\end{proof}

\subsection{Uniform covariance and arithmetic separation of scales}
\label{subsec:pf-covariance}

\begin{lemma}[Uniform long-window moments]
\label{lem:pf-moments}
Let $F:\mathbb R\to\mathbb R^D$ be the $S^2$ limit of finite
exponential sums with no constant frequency, and write its grouped
coefficients as $z_\gamma$, with
$z_{-\gamma}=\overline{z_\gamma}$ and $\sum_\gamma\|z_\gamma\|^2<\infty$.
Then its long-window mean is zero and its long-window second moment is
$\sum_\gamma z_\gamma z_\gamma^*$, uniformly over the window origin.
The same limits hold on a half-line for $F+R$ if $R$ belongs to $L^2$
on that half-line.
\end{lemma}

\begin{proof}
For a finite sum, the assertion follows by expanding products and using
\[
 \left|\frac1L\int_u^{u+L}e^{i\omega t}\,dt\right|
                 \leq\frac2{L|\omega|}\qquad(\omega\ne0),
\]
independently of $u$.  If $P$ approximates $F$ in $S^2$, covering an
interval of length $L\geq1$ by at most $L+1$ unit intervals gives
\[
 \frac1L\int_u^{u+L}\|F-P\|^2\,dt\leq2\|F-P\|_{S^2}^2.
\]
Cauchy--Schwarz controls the errors in the first and second moments;
the local second-moment bounds of $F$ and $P$ stay bounded during the
approximation.  Pass first to large $L$ and then to the $S^2$ limit.
For the last assertion,
$L^{-1}\int_u^{u+L}\|R\|^2\leq L^{-1}\|R\|_{L^2}^2$ uniformly over
allowed origins, and the same Cauchy--Schwarz argument applies.
\end{proof}

Lemmas~\ref{lem:pf-bins} and \ref{lem:pf-transfer} apply to the vector
in \eqref{eq:pf-deletion}.  Its Fourier coefficients are exactly
$z(\gamma)$ for $\gamma\notin J\cup\{0\}$.  The decay of the bin
masses proves square summability of these coefficients.  Lemma
\ref{lem:pf-moments} therefore establishes \eqref{eq:pf-moments} and
\eqref{eq:pf-cov}, including their uniformity.  Pairing opposite
frequencies makes $C_J$ real symmetric.  Positive \emph{definiteness}
requires arithmetic input, which we give next.

For a fixed real $u\geq2$, the fixed-argument Landau formula under
Dirichlet GRH is
\begin{equation}
\label{eq:pf-landau}
 \sum_{T<\gamma\leq T+H}m_\chi(\gamma)e^{i\gamma\log u}
 =-\frac{H}{2\pi\sqrt{u}}\Lambda(u)\chi^*(u)
                       +O_{q,u}(\log(2T)),
 \qquad 1\leq H\leq T.
\end{equation}
Here $\Lambda(u)\chi^*(u)$ is zero unless $u$ is an integral prime
power.  This is the fixed-argument specialization of the
Landau--Gonek formula for Dirichlet $L$-functions
\cite[Theorem~1.2]{Banks2023}; it is also the input to the article's
multiscale coercivity proof \cite{ProveItArticle}.  At a fixed
non-prime-power argument the distance to the neighboring prime powers
is positive.  Endpoints change the sum by at most $O_q(\log(2T))$, by
\eqref{eq:pf-localcount}.

\begin{lemma}[Prime probes]
\label{lem:pf-probes}
For each $a\in G$ one can choose a prime $p_a\equiv a\pmod q$ such
that all numbers $p_a e^{h_k-h_j}$ are at least $2$, and are integral
prime powers only when $j=k$.
\end{lemma}

\begin{proof}
For fixed $R>0$, $R\ne1$, only finitely many primes $p$ make $Rp$ an
integral prime power.  If $R$ is irrational there are none.  Otherwise
write $R=b/c$ in lowest terms.  If $c>1$, integrality forces $c\mid p$
and leaves at most one prime.  If $c=1$, then $b\geq2$ and $bp$ can
be a prime power only if $p$ divides $b$.  There are finitely many
off-diagonal ratios $e^{h_k-h_j}\ne1$.  Exclude their exceptional
primes, take $p$ large enough for the lower bound $2$, and use
Dirichlet's theorem in the reduced class $a$.
\end{proof}

\begin{lemma}[An unsmoothed band is coercive]
\label{lem:pf-coercivity}
There are constants $c,K,T_0>0$, depending only on $q$ and the shifts,
such that, for $T\geq T_0$ and $K\log(2T)\leq H\leq T$,
\begin{equation}
\label{eq:pf-coercivity}
 2\sum_{\gamma\in\Gamma_q\cap(T,T+H]}
 \frac{|A_w(\gamma)|^2}{\tfrac14+\gamma^2}
       \geq\frac{cH}{T^2\log(2T)}\|w\|_2^2,
 \quad
 A_w(\gamma)=\sum_{j,a}w_{j,a}e^{i\gamma h_j}v_a(\gamma).
\end{equation}
This holds for complex coefficient arrays $w$ as well as real ones.
\end{lemma}

\begin{proof}
Let $B=\Gamma_q\cap(T,T+H]$, and use the primes from
Lemma~\ref{lem:pf-probes}.  Expanding \eqref{eq:pf-v} and applying
\eqref{eq:pf-landau} at the finitely many fixed arguments
$p_a e^{h_k-h_j}$ gives
\begin{align}
\label{eq:pf-recover}
 \sum_{\gamma\in B}A_w(\gamma)
              e^{i\gamma(\log p_a-h_j)}
 &=-\frac{H\log p_a}{2\pi\sqrt{p_a}}w_{j,a}
                    +O_{q,h}\bigl(\|w\|_2\log(2T)\bigr).
\end{align}
All off-diagonal main terms vanish.  On the diagonal,
$\chi^*(p_a)=\chi(a)$ because $p_a$ is prime to $q$, and character
orthogonality isolates $w_{j,a}$.

Choose a coordinate with $|w_{j,a}|\geq\|w\|_2/\sqrt D$.
The positive number
$\min_a\log p_a/(2\pi\sqrt{p_a})$ is fixed.  By taking $K$ large
enough, the absolute value on the left of \eqref{eq:pf-recover} is
at least $c_1H\|w\|_2$.  The local zero count, summed over the band,
gives $|B|\leq c_2H\log(2T)$.  Cauchy--Schwarz over its
\emph{distinct} ordinates therefore proves
\[
 \sum_{\gamma\in B}|A_w(\gamma)|^2
                 \geq c_3\frac{H}{\log(2T)}\|w\|_2^2.
\]
Since $\tfrac14+\gamma^2\leq5T^2$ in this band for $T\geq1$,
division by the denominator proves \eqref{eq:pf-coercivity}.
\end{proof}

For real $w$, pairing $\gamma$ and $-\gamma$ in
\eqref{eq:pf-cov} shows that $w^{\mathsf T}C_Jw$ is at least the
left side of \eqref{eq:pf-coercivity} for every band above the deleted
frequencies.  This proves \eqref{eq:pf-band} and positive
definiteness, including after arbitrary fixed finite deletion.  The
argument separates scales by prime powers; it does not separate
individual nearby zeros.

\subsection{Two-sided occupation and failure of exact linear prediction}
\label{subsec:pf-occupation}

We complete the proof of Theorem~\ref{thm:pf-main} by a moment
argument which does not require pointwise boundedness of the
unsmoothed errors.

\begin{lemma}[Occupation from a fourth-moment bound]
\label{lem:pf-occupation}
Suppose a real vector-valued function $Y$ has uniform long-window mean
zero and a uniform long-window second moment $C$ with
$\lambda=\lambda_{\min}(C)>0$.  Suppose also that, for every real
unit $w$ and every sufficiently long allowed interval $I$,
$|I|^{-1}\int_I|w\cdot Y|^4\leq M^4$, where $M\geq1$ is independent
of $w$ and $I$.  Then positive and negative directional occupation
hold uniformly, with permissible constants
\begin{equation}
\label{eq:pf-constants}
 a=\frac{(\lambda/2)^{3/2}}{M^2},
 \qquad \varepsilon=\frac a8,
 \qquad \eta=\left(\frac{a}{8M}\right)^2.
\end{equation}
\end{lemma}

\begin{proof}
Let $\mathbb E$ denote average over the interval and set $f=w\cdot Y$.
Uniformity of the first two moments allows us to make
$\mathbb Ef^2\geq\lambda/2$ and $|\mathbb Ef|\leq a/2$
simultaneously for all unit $w$.  Interpolation between $L^1$ and
$L^4$, equivalently H\"older's inequality, gives
\[
 \mathbb E|f|\geq
       \frac{(\mathbb Ef^2)^{3/2}}{(\mathbb Ef^4)^{1/2}}
       \geq a.
\]
Writing $f_+=\max(f,0)$ and $f_-=\max(-f,0)$, we obtain
$\mathbb Ef_+\geq a/4$ and $\mathbb Ef_-\geq a/4$.
With $\varepsilon=a/8$ this implies
\[
 \frac a8\leq\mathbb E\bigl[f_+\mathbf1_{\{f\geq\varepsilon\}}\bigr]
       \leq(\mathbb Ef^2)^{1/2}
                       \mathbb P(f\geq\varepsilon)^{1/2}
       \leq M\mathbb P(f\geq\varepsilon)^{1/2}.
\]
Here $\mathbb P$ is relative Lebesgue measure on the interval.
The same calculation for $-f$ proves both assertions with the
specified $\eta$.
\end{proof}

The Fourier vector obtained from \eqref{eq:pf-F} by shifts and finite
deletion has bounded $S^4$ norm.  The remainder from
Lemma~\ref{lem:pf-transfer} has bounded local fourth moments too.
Covering a long interval by unit intervals therefore supplies a
common $M$ for $Y=\mathcal X_J$.  The covariance and mean were
established uniformly above, so Lemma~\ref{lem:pf-occupation} proves
\eqref{eq:pf-positive}--\eqref{eq:pf-negative}.  If the closed
convex hull of the essential values failed to contain the
$\varepsilon$-ball, separation would give a unit direction with
essential supremum less than $\varepsilon$, contradicting
\eqref{eq:pf-positive}.  This completes the main theorem.

\begin{corollary}[No exact finite linear explanation]
\label{cor:pf-prediction}
For a nonzero real array $w=(w_{j,a})$, an arbitrary real finite
trigonometric polynomial $P$, and an arbitrary real constant $b$,
the long-window mean square of
\[
 \sum_{j,a}w_{j,a}X_a(t+h_j)-P(t)-b
\]
has a strictly positive limit, uniformly over the allowed origins.
In particular, no nonzero contrast at one fixed scale can be
predicted exactly in this mean-square sense by real linear
combinations of the other finitely many scales and a finite
trigonometric correction.
\end{corollary}

\begin{proof}
Let $J$ contain all nonzero frequencies of $P$ which belong to
$\Gamma_q$, together with their negatives.  Frequencies outside
this finite set retain precisely the coefficients of
$w\cdot\mathcal X_J$.  The uniform Fourier mean-square calculation
of Lemma~\ref{lem:pf-moments}, including any extra frequencies of
$P$, shows that the residual mean square is at least
$w^{\mathsf T}C_Jw>0$.  For a predictor, collect all coefficients in
one array; its block at the predicted scale is the prescribed
nonzero contrast, so the full array cannot vanish.
\end{proof}

This corollary concerns every \emph{fixed} finite correction.  It
does not give a positive lower bound uniform over all corrections:
successively larger Fourier truncations do approximate the error
in mean square.  Likewise, the occupation theorem describes
individual linear contrasts of centered counts.  Positive
covariance and two-sided directional occupation alone do not force
the limiting values to fill an open set, and therefore do not
imply inclusivity of every ordinary prime race.

\section{An exponential information gap for wide symplectic blocks}
\label{sec:symplectic-information}

The second result concerns a finite linear-algebra mechanism developed in the
source project.  A rectangular block of a uniformly random symplectic
transformation can have almost the same rank distribution as a matrix of
independent uniform entries, while the two complete matrix distributions are
almost mutually singular.  We determine both discrepancies and the number of
rank observations needed to distinguish the models.

The exact arbitrary-frame compression law and a square-block version of rank
blindness already appear in the supplied article
\cite[sections on symplectic compressions and Lagrangian projections]{ProveItArticle}.
The square result gives a positive limiting discrepancy between complete
matrices.  Here the number of columns exceeds the number of rows by a growing
amount.  In this regime the complete-matrix discrepancy tends to one, and we
prove a sharp rank-only sample scale.  The finite orbit count below is a
specialization of the article's law, for which we give a direct proof.  The
uniform asymptotic separation and statistical sample scale are the additional
conclusions.  They are unconditional finite-field theorems; they do not use
Dirichlet GRH.  Classical uniform-matrix rank distributions provide the
comparison law \cite{FulmanGoldstein2015}; prescribed-form sampling is also
studied in \cite{Ishizuka2026}.

\subsection{The two random-matrix models}

Fix a prime power $q$, integers
\begin{equation}
  n\ge c\ge1,\qquad b=n+c\text{ even},\qquad m=b/2,
  \qquad Q=q^{-1},\qquad h=n-c.
  \label{symp:eq:parameters}
\end{equation}
In particular $h$ is even and $0\le h\le n-1$.  Let $H$ be a
nondegenerate alternating $b\times b$ matrix over $\mathbb F_q$.  Here
``alternating'' means $x^{\mathsf T}Hx=0$ for every $x$; in characteristic two
this includes the zero-diagonal condition.  Set
\[
 J_n=\begin{pmatrix}0&I_n\\-I_n&0\end{pmatrix}.
\]
Choose a $2n\times b$ matrix $X_0$ satisfying
$X_0^{\mathsf T}J_nX_0=H$.  Such a matrix exists since $b\le2n$; its columns
are independent because $H$ is nondegenerate.  For uniform
$g\in\operatorname{Sp}_{2n}(\mathbb F_q)$, write
\begin{equation}
 gX_0=\begin{pmatrix}A\\B\end{pmatrix},
 \qquad A,B\in\operatorname{Mat}_{n,b}(\mathbb F_q).
 \label{symp:eq:model}
\end{equation}
Let $U$ be uniform on $\operatorname{Mat}_{n,b}(\mathbb F_q)$, so its entries
are independent and uniform.  The two laws are independent of the particular
choice of $X_0$ with the prescribed Gram matrix.

For probability measures on a finite set we use
\[
 \operatorname{TV}(P,V)=\frac12\sum_x|P(x)-V(x)|,
 \qquad
 \chi^2(P\Vert V)=\sum_x\frac{(P(x)-V(x))^2}{V(x)}
\]
when $V$ is positive on the support of $P$.  Write
$\operatorname{TV}(A,U)$ for the distance between their laws.  Define
\begin{equation}
 a_j=\prod_{i=1}^j(1-Q^i),\quad a_0=1,\quad
 a_\infty=\prod_{i=1}^\infty(1-Q^i)>0,\quad
 C=\prod_{j=n-m+1}^{n}(1-Q^{2j}).
 \label{symp:eq:products}
\end{equation}
All finite products with empty index sets are one.  For a proposed matrix
$M$, let $K=\ker M$ and $k=\dim K$.  Necessarily $c\le k\le b$.
We call $K$ \emph{$H$-isotropic} if $x^{\mathsf T}Hy=0$ for every $x,y\in K$.

\begin{theorem}[Exact law and optimal distinguishing events]
\label{symp:thm:finite}
In the model \eqref{symp:eq:model}, the likelihood relative to a uniform
matrix is
\begin{equation}
 \frac{\Pr(A=M)}{q^{-nb}}
 =\mathbf 1_{\{\ker M\text{ is }H\text{-isotropic}\}}
   \frac{q^{\binom{k}{2}}}{C}.
 \label{symp:eq:density}
\end{equation}
Consequently,
\begin{equation}
 \operatorname{TV}(A,U)
 =\Pr\bigl(H|_{\ker U\times\ker U}\ne0\bigr).
 \label{symp:eq:fulltv}
\end{equation}
For $c\le k\le b$ put
\begin{equation}
 u(k)=Q^{k(k-c)}\frac{a_na_b}{a_ka_{k-c}a_{b-k}},\qquad
 R_{b,k}=\prod_{i=0}^{k-1}
          \frac{1-Q^{b-2i}}{1-Q^{b-i}}\quad(k\le m),
 \label{symp:eq:rank-formulas}
\end{equation}
and set $R_{b,k}=0$ for $k>m$.  Then
\begin{equation}
 \Pr(\dim\ker U=k)=u(k),\qquad
 \Pr(\dim\ker A=k)=u(k)\frac{R_{b,k}}C.
 \label{symp:eq:rank-law}
\end{equation}
Full row rank is the only rank class with more probability under $A$ than
under $U$.  In particular, with
\[
 D=\prod_{j=m+1}^n(1-Q^{2j}),\qquad
 E=\prod_{j=n+1}^{n+c}(1-Q^j),
\]
the exact rank discrepancy is
\begin{equation}
 \Delta:=\operatorname{TV}(\operatorname{rank}A,\operatorname{rank}U)
 =\frac{a_n}{a_c}\bigl(D^{-1}-E\bigr).
 \label{symp:eq:rank-tv}
\end{equation}
The overlap of the complete matrix laws is
\begin{equation}
 \Omega:=1-\operatorname{TV}(A,U)
  =\sum_{k=c}^{m}u(k)Q^{\binom{k}{2}}R_{b,k}.
 \label{symp:eq:overlap}
\end{equation}
\end{theorem}

\subsection{Counting frames and fibres}

We first prove the likelihood without using the general compression formula.
The symplectic group acts transitively on the frames $X$ satisfying
$X^{\mathsf T}J_nX=H$: identify the nondegenerate spans of two frames by the
map taking corresponding columns to each other, and extend that isometry
over their symplectic orthogonal complements.  A fixed frame has stabilizer
$\operatorname{Sp}_{2(n-m)}(\mathbb F_q)$, acting on its complement.  The
standard count
\[
 |\operatorname{Sp}_{2t}(\mathbb F_q)|
   =q^{t^2}\prod_{j=1}^{t}(q^{2j}-1)
\]
therefore gives
\begin{equation}
 \#\{X:X^{\mathsf T}J_nX=H\}
 =q^{2nb-b(b-1)/2}C.
 \label{symp:eq:frame-count}
\end{equation}
The group-order formula itself follows by successively choosing a nonzero
first vector, a partner pairing to one with it, and a symplectic basis in
their orthogonal complement.

Fix $A=M$ of rank $r=b-k$.  The equation for the remaining rows is
\begin{equation}
 M^{\mathsf T}B-B^{\mathsf T}M=H.
 \label{symp:eq:gram-fibre}
\end{equation}
Its restriction to $K\times K$ is zero, so $H$-isotropy of $K$ is necessary.
Choose invertible matrices $P,T$ with
$PMT=\left(\begin{smallmatrix}I_r&0\\0&0\end{smallmatrix}\right)$.
Replacing $B$ by $P^{-\mathsf T}BT$ and $H$ by $T^{\mathsf T}HT$ preserves
\eqref{symp:eq:gram-fibre}.  Partition rows into sizes $r,n-r$ and columns
into sizes $r,k$.  In these coordinates the equation prescribes
\[
 B_{11}-B_{11}^{\mathsf T}=H_{11},\qquad B_{12}=H_{12},
 \qquad H_{22}=0.
\]
If $K$ is isotropic, the last equation holds.  The first has
$q^{r(r+1)/2}$ solutions, since the diagonal and one entry in each
off-diagonal pair are free.  The last $n-r$ rows of $B$ are entirely free.
Thus the number of solutions is
\begin{equation}
 q^{r(r+1)/2+(n-r)b}
  =q^{nb-b(b-1)/2+\binom{k}{2}}.
 \label{symp:eq:fibre-count}
\end{equation}
Every such solution gives an independent frame: a vector in the kernel of
$\left(\begin{smallmatrix}M\\B\end{smallmatrix}\right)$ belongs to
$\operatorname{rad}H=0$.  This argument requires no division by two and
applies in characteristic two.  Dividing \eqref{symp:eq:fibre-count} by
\eqref{symp:eq:frame-count} proves \eqref{symp:eq:density}.

Since $C<1$, every positive likelihood in \eqref{symp:eq:density} is
strictly larger than one.  The entire deficit relative to the uniform law
therefore lies outside the support, proving \eqref{symp:eq:fulltv}.  If $Z$
is a basis matrix for $\ker M$, the optimal event is simply
$Z^{\mathsf T}HZ\ne0$.  It can be evaluated using a polynomial number of
finite-field operations.

\subsection{The rank law and its sign pattern}

To count uniform matrices with kernel dimension $k$, choose a $k$-subspace
of $\mathbb F_q^b$ and an injective map from its quotient to
$\mathbb F_q^n$.  Dividing the resulting count by $q^{nb}$ gives $u(k)$
in \eqref{symp:eq:rank-formulas}.  Conditional on $k$, the kernel is a
uniform $k$-subspace.

For an independent isotropic $k$-frame in the nondegenerate $b$-space,
after $i$ vectors have been chosen the next vector must lie in their common
orthogonal complement and outside their span.  There are
$q^{b-i}-q^i$ choices.  Comparing with all independent $k$-frames gives
\begin{equation}
 \Pr(K\text{ is }H\text{-isotropic}\mid\dim K=k)
 =\frac{\prod_{i=0}^{k-1}(q^{b-i}-q^i)}
        {\prod_{i=0}^{k-1}(q^b-q^i)}
 =Q^{\binom{k}{2}}R_{b,k}
 \label{symp:eq:isotropy}
\end{equation}
for $k\le m$, and the probability is zero for $k>m$.  Averaging
\eqref{symp:eq:density} proves \eqref{symp:eq:rank-law}; summing
\eqref{symp:eq:isotropy} proves \eqref{symp:eq:overlap}.

It remains to identify the sign of every rank discrepancy.  Cancellation
of factors gives
\begin{equation}
 \frac{R_{b,c}}C=(DE)^{-1}>1.
 \label{symp:eq:fullrank-ratio}
\end{equation}
For $1\le k<m$ the successive ratio
\[
 \frac{R_{b,k+1}}{R_{b,k}}
  =\frac{1-Q^{b-2k}}{1-Q^{b-k}}
\]
is strictly smaller than one.  We prove that $R_{b,c+1}<C$, after which
all remaining inequalities follow.  If $c=n$, then $c=m$ and this is
immediate.  Otherwise $h=n-c\ge2$, and direct cancellation gives
\begin{equation}
 \frac{R_{b,c+1}}C
 =\frac{1-Q^{n-c}}
 {\displaystyle\prod_{j=n}^{n+c}(1-Q^j)
  \displaystyle\prod_{j=m+1}^{n}(1-Q^{2j})}.
 \label{symp:eq:first-deficiency}
\end{equation}
The sum of the quantities subtracted from one in this denominator, divided
by $Q^{n-c}$, is at most
\begin{equation}
 \frac{Q^c(1-Q^{c+1})}{1-Q}
 +\frac{Q^{2c+2}}{1-Q^2}<1.
 \label{symp:eq:sign-bound}
\end{equation}
Indeed, for $c=1$ the upper bound is at most
$3/4+1/12=5/6$; for $c\ge2$ it is at most
$1/2+1/48=25/48$.  We used only $Q\le1/2$.  The inequality
$\prod_i(1-x_i)\ge1-\sum_i x_i$ now shows that the denominator in
\eqref{symp:eq:first-deficiency} exceeds its numerator.  Hence full row
rank is the sole rank class with positive discrepancy.  Multiplying its
likelihood excess by $u(c)=a_b/a_c$ proves
\eqref{symp:eq:rank-tv}, completing the proof of
Theorem~\ref{symp:thm:finite}.

\subsection{Uniform separation as the excess width grows}

\begin{theorem}[Almost singular matrices with almost identical ranks]
\label{symp:thm:separation}
Fix $q$.  As $c\to\infty$, uniformly for every $n\ge c$ with $n+c$ even,
\begin{equation}
 \Delta=\frac{q^{-n}}{q-1}\bigl(1+O_q(q^{-c})\bigr).
 \label{symp:eq:rank-asymptotic}
\end{equation}
An explicit estimate, valid throughout \eqref{symp:eq:parameters}, is
\begin{equation}
 \left|(q-1)q^n\Delta-1\right|\le3q^{-c}.
 \label{symp:eq:rank-explicit}
\end{equation}
Define
\begin{align}
 \Omega_0&=q^{-\binom c2}\frac{a_n}{a_c}
  \prod_{j=(n-c)/2+1}^{(n+c)/2}(1-q^{-2j}),
  \label{symp:eq:leading-overlap}\\
 \eta_q(c)&=\frac{a_\infty^{-2}q^{-(2c+1)}}{1-q^{-(2c+4)}}.
  \label{symp:eq:overlap-error}
\end{align}
Then
\begin{equation}
 \Omega_0\le\Omega\le\Omega_0\bigl(1+\eta_q(c)\bigr).
 \label{symp:eq:overlap-certified}
\end{equation}
In particular, for the even integer $h=n-c$ put
\[
 B_{q,h}=\prod_{j=h/2+1}^{\infty}(1-q^{-2j}).
\]
Uniformly for $h\ge0$,
\begin{equation}
 1-\operatorname{TV}(A,U)
 =B_{q,h}\,q^{-\binom c2}\bigl(1+O_q(q^{-c})\bigr).
 \label{symp:eq:full-asymptotic}
\end{equation}
Thus $\operatorname{TV}(\operatorname{rank}A,\operatorname{rank}U)\to0$
and $\operatorname{TV}(A,U)\to1$.  If also $n-c\to\infty$, the leading
factor $B_{q,h}$ tends to one.  No such additional hypothesis is needed for
the separation itself.
\end{theorem}

\begin{proof}
For the rank estimate, write
\[
 S_E=\sum_{j=n+1}^{n+c}Q^j
     =\frac{Q^n(1-Q^c)}{q-1},\qquad
 S_D=\sum_{j=m+1}^{n}Q^{2j}
     \le\frac{Q^{n+c}}{q^2-1}.
\]
The elementary product estimates give
\begin{gather}
 S_E-\tfrac12S_E^2\le1-E\le S_E,\qquad
 0\le D^{-1}-1\le\frac{S_D}{1-S_D},
 \label{symp:eq:product-bounds}\\
 1-\frac{Q^c}{q-1}\le\frac{a_n}{a_c}\le1.
 \nonumber
\end{gather}
For the first lower bound, the sum of pairwise products is at most
$S_E^2/2$, and the two-term Bonferroni inequality applies.  Since
$n+c\ge2$, we have $S_D\le1/12$.  Divide
\eqref{symp:eq:rank-tv} by $Q^n/(q-1)$.  The upper bound from
\eqref{symp:eq:product-bounds} is at most
$1+(4/11)Q^c$.  The lower bound is at least
\[
 \left(1-\frac{Q^c}{q-1}\right)
 \left(1-Q^c-\frac{Q^n}{2(q-1)}\right)
 \ge1-\frac52Q^c.
\]
The weaker symmetric bound \eqref{symp:eq:rank-explicit} follows.

The $k=c$ summand of \eqref{symp:eq:overlap} is exactly
$\Omega_0$.  To bound the others, put $k=c+j$.  Since $R_{b,c+j}\le R_{b,c}$,
and
\begin{equation}
 \frac{u(c+j)}{u(c)}
  =Q^{j(c+j)}\frac{a_ca_n}{a_{c+j}a_ja_{n-j}}
  \le a_\infty^{-2}Q^{j(c+j)},
 \label{symp:eq:overlap-ratio}
\end{equation}
the ratio of the $j$th summand to $\Omega_0$ is at most
\begin{equation}
 a_\infty^{-2}Q^{2cj+(3j^2-j)/2}.
 \label{symp:eq:overlap-tail}
\end{equation}
For $j\ge1$, successive exponents increase by
$2c+3j+1\ge2c+4$.  A geometric majorant proves
\eqref{symp:eq:overlap-certified}.  Finally,
$a_n/a_c=1+O_q(Q^c)$, while the finite even product in
\eqref{symp:eq:leading-overlap} differs from $B_{q,h}$ by relative
$O_q(Q^{n+c})$.  Because $n+c\ge2c$, these estimates prove
\eqref{symp:eq:full-asymptotic}, uniformly in $h$.
\end{proof}

The boundary $c=n$ is included.  Here the source frame fills the ambient
space, so $A$ has full row rank surely.  The exact formulas reduce to
\begin{equation}
 \Delta=1-\prod_{j=n+1}^{2n}(1-q^{-j}),\qquad
 \Omega=q^{-\binom n2}\prod_{j=1}^{n}(1-q^{-2j}).
 \label{symp:eq:saturation}
\end{equation}
The second equality also counts the uniformly random matrices whose
$n$-dimensional kernels are Lagrangian.  More generally, when $h$ remains
fixed, the factor $B_{q,h}$ records the finite distance to this boundary.

\subsection{Sharp sample complexity when only ranks are observed}

Small one-observation total variation by itself does not identify the right
sample scale.  The following calculation does so.

\begin{theorem}[Rank-only information and its sharp sample scale]
\label{symp:thm:sample}
Let $P$ be the law of $\operatorname{rank}A$ and let $V$ be the law of
$\operatorname{rank}U$.  For fixed $q$ and $c\to\infty$, uniformly over
\eqref{symp:eq:parameters},
\begin{equation}
 \chi^2(P\Vert V)
  =\frac{q^{-(2n-c)}}{q-1}\bigl(1+O_q(q^{-c})\bigr).
 \label{symp:eq:chi-square}
\end{equation}
For $t$ independent matrix observations, with only their ranks retained,
the following conclusions hold.
\begin{enumerate}
\item If $t=o(q^{2n-c})$, the total variation distance between the two
 joint rank laws tends to zero.  Under equal prior probabilities, every
 rank-only decision rule consequently has error probability tending to
 $1/2$ or larger.
\item If $t/q^{2n-c}\to\infty$, a threshold on the number of matrices
 failing to have full row rank distinguishes the models with error
 probability tending to zero.
\end{enumerate}
For any fixed target error strictly between zero and $1/2$, the required
number of rank observations is therefore $\Theta_q(q^{2n-c})$, with the
implicit constants also allowed to depend on the target error.
By contrast, one complete matrix has optimal equal-prior error
\begin{equation}
 \frac{\Omega}{2}
 =\frac{B_{q,n-c}}2\,q^{-\binom c2}
      \bigl(1+O_q(q^{-c})\bigr).
 \label{symp:eq:one-matrix-risk}
\end{equation}
\end{theorem}

\begin{proof}
Use row deficiency $j$ instead of rank, so that the kernel has dimension
$k=c+j$.  Write
\[
 u_j=u(c+j),\qquad r_j=\frac{R_{b,c+j}}C.
\]
Then $P(j)=u_jr_j$, $V(j)=u_j$, and
\begin{equation}
 \chi^2(P\Vert V)=\sum_{j=0}^{n}u_j(r_j-1)^2.
 \label{symp:eq:chi-sum}
\end{equation}
We show that the term $j=1$ supplies the leading contribution.

First, the explicit uniform rank law gives
\begin{equation}
 u_1=\frac{Q^c}{q-1}(1-Q^n)
         \prod_{i=c+2}^{n+c}(1-Q^i)
     =\frac{Q^c}{q-1}\bigl(1+O_q(Q^c)\bigr).
 \label{symp:eq:u-one}
\end{equation}
From \eqref{symp:eq:fullrank-ratio}, $r_0=1+O_q(Q^n)$.  The successive
ratio formula gives, including when $h=0$,
\begin{equation}
 r_1=r_0\frac{1-Q^h}{1-Q^n},\qquad
 1-r_1=Q^h\bigl(1+O_q(Q^c)\bigr).
 \label{symp:eq:r-one}
\end{equation}
For clarity, the second estimate follows by writing
$r_0/(1-Q^n)=1+O_q(Q^n)$: subtracting its product with $1-Q^h$ from one
leaves $Q^h+O_q(Q^n)$, and $Q^n/Q^h=Q^c$.  At $h=0$ the identity
$r_1=0$ makes the conclusion exact.  Equations
\eqref{symp:eq:u-one} and \eqref{symp:eq:r-one} yield
\begin{equation}
 u_1(1-r_1)^2
 =\frac{Q^{c+2h}}{q-1}\bigl(1+O_q(Q^c)\bigr).
 \label{symp:eq:chi-leading}
\end{equation}
The $j=0$ contribution is $O_q(Q^{2n})$, a relative
$O_q(Q^c)$ error.

For the other terms we need an estimate that remains valid near the support
boundary.  For every $j\ge1$,
\begin{equation}
 0\le1-r_j\le\frac{Q^{h-2(j-1)}}{1-Q^2}.
 \label{symp:eq:likelihood-defect}
\end{equation}
If $c+j\le m$, then $r_j\ge R_{b,c+j}$, and removing the denominators in
the product for $R$ gives
\[
 R_{b,c+j}\ge\prod_{i=0}^{c+j-1}(1-Q^{b-2i}).
\]
Its defect from one is at most the sum of the displayed powers of $Q$,
which is bounded by the right side of
\eqref{symp:eq:likelihood-defect}.  If $c+j>m$, then $r_j=0$ and
$h-2(j-1)\le0$, so the same bound still holds.  The uniform rank formula
also gives
\begin{equation}
 u_j=Q^{j(c+j)}
       \frac{a_n}{a_{n-j}}\frac{a_b}{a_{c+j}}\frac1{a_j}
     \le a_\infty^{-1}Q^{j(c+j)}.
 \label{symp:eq:rank-tail-bound}
\end{equation}
Consequently, for $j\ge2$, the summand in
\eqref{symp:eq:chi-sum}, divided by $Q^{c+2h}$, is at most
\[
 \frac{Q^{(j-1)c+(j-2)^2}}{a_\infty(1-Q^2)^2}.
\]
Their sum is bounded by
\begin{equation}
 \frac{1}{a_\infty(1-Q^2)^2}
       \frac{Q^c}{1-Q^{c+1}},
 \label{symp:eq:chi-remainder}
\end{equation}
since successive exponents increase by at least $c+1$.  Together with
\eqref{symp:eq:chi-leading}, and $c+2h=2n-c$, this proves
\eqref{symp:eq:chi-square} with uniform error.

For the impossibility assertion, independence gives the exact identity
\[
 1+\chi^2(P^{\otimes t}\Vert V^{\otimes t})
  =\bigl(1+\chi^2(P\Vert V)\bigr)^t.
\]
By Cauchy--Schwarz,
\[
 \operatorname{TV}(P^{\otimes t},V^{\otimes t})
 \le\frac12\sqrt{\bigl(1+\chi^2(P\Vert V)\bigr)^t-1},
\]
which tends to zero when $t=o(q^{2n-c})$.  Optimal equal-prior error is
$(1-\operatorname{TV})/2$, proving the first assertion.

For achievability, let
\[
 p_U=\Pr(\operatorname{rank}U<n),\qquad
 p_A=\Pr(\operatorname{rank}A<n).
\]
The exact sign pattern gives $p_U-p_A=\Delta$.  The uniform full-row-rank
product and a union bound show
\[
 0\le p_A\le p_U
   =1-\prod_{j=c+1}^{n+c}(1-Q^j)
   \le\frac{Q^c}{q-1}.
\]
From $t$ observations compute the fraction failing to have full row rank,
and compare it with $(p_A+p_U)/2$.  Under either hypothesis its variance is
at most $p_U/t$.  Chebyshev's inequality bounds each error probability by
\begin{equation}
 \frac{4p_U}{t\Delta^2}
   =O_q\left(\frac{q^{2n-c}}t\right).
 \label{symp:eq:chebyshev-test}
\end{equation}
This tends to zero in the claimed range.  A sufficiently small constant
multiple of the scale in the chi-square bound, and a sufficiently large
constant multiple in \eqref{symp:eq:chebyshev-test}, give the matching
fixed-error lower and upper bounds.

Finally, the complete-matrix test selects the symplectic model precisely
when the observed kernel is $H$-isotropic.  It never errs under that model,
and errs with probability $\Omega$ under the uniform model.  It is optimal
by \eqref{symp:eq:density}, so its equal-prior error is $\Omega/2$.
Theorem~\ref{symp:thm:separation} gives
\eqref{symp:eq:one-matrix-risk}.
\end{proof}

\subsection{Interpretation and relation to the source}

For example, take any admissible sequence with $c/n\to\alpha\in(0,1)$.
The one-observation rank discrepancy is asymptotic to
$q^{-n}/(q-1)$, and the rank-only sample scale is
$q^{(2-\alpha)n+o(n)}$.  Yet the optimal error from a single complete matrix
is
\[
 \frac12q^{-\alpha^2n^2/2+o(n^2)}.
\]
The information resides in pairings between vectors in the kernel, rather
than in its dimension.  Kernel dimension is almost identically distributed
in the two models; the restriction of the known alternating form to that
kernel separates them.

The supplied article establishes the square case and the general exact
compression law.  The present family keeps the target Lagrangian, takes a
nondegenerate source, and allows a growing excess of columns.  This yields
an explicit solution of one growing-dimension regime, including its
saturation boundary and its rank-only statistical complexity.  It does not
determine the limiting law of arbitrary pairs of growing, possibly
degenerate, restricted forms.  The literature comparison likewise supports
a new extension relative to the supplied corpus; it is not a certification
of worldwide priority.

\section{Exact computation and reproducibility}\label{sec:checks}
The accompanying program \texttt{code/verify\_symplectic.py} uses only the
Python standard library. It performs two independent types of verification.
First, it recursively enumerates actual ordered symplectic pairs in successive
orthogonal complements, projects their coordinates, and obtains the empirical
law as an exact count. Separately, it row-reduces every possible output matrix,
constructs its kernel, and tests its isotropy. The complete probability of
each matrix is compared with \eqref{symp:eq:density}; aggregation checks the
rank law, total variation, and chi-square divergence.

The nine enumeration cases use these parameters:
\[
\begin{array}{c|l}
q & (n,b)\\\hline
2 & (1,2),\ (2,2),\ (2,4),\ (3,4)\\
3 & (1,2),\ (2,2),\ (2,4)\\
4 & (1,2),\ (2,2)
\end{array}
\]
The square cases are auxiliary checks of the counting identity.
In total, \textbf{313,170 prescribed-Gram
frames and 11,295 individual matrices} were examined, with no mismatches.
The quaternary cases use actual extension-field arithmetic in
$\mathbb F_4=\mathbb F_2[x]/(x^2+x+1)$, not arithmetic modulo four.

Second, exact rational arithmetic checks normalization, the strict rank sign
threshold, and the closed expression \eqref{symp:eq:rank-tv} for all 288
admissible triples with $q\in\{2,3,4,5\}$, $1\le n\le16$, $1\le c\le n$,
and $n+c$ even. Additional sequences compare the asymptotics with exact
fractions, including $n=c$. These are implementation checks; the proofs in
Section~\ref{sec:symplectic-information} establish the all-parameter claims.

\begin{table}[htbp]
\centering\small
\begin{tabular}{rrlll}
\toprule
$n$ & $c$ & Rank TV $\Delta$ & Full overlap $\Omega$ & Rank $\chi^2$\\
\midrule
8 & 4 & $3.519397\times10^{-3}$ & $1.447624\times10^{-2}$ & $2.508953\times10^{-4}$ \\
12 & 6 & $2.378620\times10^{-4}$ & $2.990148\times10^{-5}$ & $3.851021\times10^{-6}$ \\
20 & 10 & $9.521233\times10^{-7}$ & $2.838477\times10^{-14}$ & $9.319256\times10^{-10}$ \\
40 & 20 & $9.094933\times10^{-13}$ & $6.372360\times10^{-58}$ & $8.673623\times10^{-19}$ \\
12 & 12 & $2.440612\times10^{-4}$ & $9.331424\times10^{-21}$ & $2.441207\times10^{-4}$ \\
\bottomrule
\end{tabular}
\caption{Binary examples from exact rational formulas, rounded for display.
The small quantity $\Omega=1-\operatorname{TV}(A,U)$ is displayed directly
so that rounding does not erase the complete-matrix distinction.}
\label{tab:numeric}
\end{table}

For $q=2,n=40,c=20$, a single complete matrix has optimal equal-prior error
$\Omega/2\approx3.18618\cdot10^{-58}$. Ranks alone have the asymptotic
sample scale $2^{60}=1\,152\,921\,504\,606\,846\,976$; this is an order of
sample complexity, not an asserted exact number for a specified finite risk.
The one-observation rank distance is only about $9.095\cdot10^{-13}$.
The explicit test computes a kernel basis $K$ and tests whether
$K^{\mathsf T}HK=0$; Gaussian elimination makes it a polynomial-time test
when the source form is known.

The archive contains the full LaTeX sources, this compiled PDF, exact
verification results, a source audit, a novelty record, and a checker for the
pinned source hashes. The arithmetic result is an analytic proof and is not
inferred from a finite sample of primes or zeros. No numerical certificate
of central orders, zero simplicity, or ordinate independence is used.


\section{Formalization route and remaining boundaries}\label{sec:formalization}
For a Lean implementation, the two proofs separate into reusable interfaces.
The arithmetic part needs the source's Dirichlet GRH theorem, the ordinary
explicit formula, a local zero-count bound, and fixed-argument Landau sums.
The new analysis can then be developed as: (i) the binned local $L^2$ estimate;
(ii) its convolution-based $L^4$ upgrade; (iii) completeness and uniform mean
limits in these norms; (iv) the Volterra estimate transferring $\theta$ to
$\pi$; and (v) the finite-dimensional occupation lemma. The prime-probe
argument supplies injectivity of each sufficiently high frequency band.
Every spectral sum groups coincident ordinates before taking a square.

The finite-field part can be formalized independently of all analytic
number theory. Its natural order is: symplectic-basis extension, the
prescribed-Gram frame count, the affine fibre of $M^{\mathsf T}B-B^{\mathsf T}M=H$,
kernel-isotropy counts, the exact probability laws, product inequalities,
and the chi-square product identity. The Python checks exercise the finite
identities with exact arithmetic; they do not replace these general proofs.

The arithmetic theorem is uniform over the location of sufficiently long
logarithmic intervals and over unit directions, with modulus, scales, and
deleted frequencies fixed. It gives no uniform constants as scales coalesce,
as the modulus grows, or as the number of scales grows. Positive definite
covariance does not imply that the limiting law has an interior, nor that all
coordinate orderings occur. The matrix theorem is uniform in $n\ge c$ for fixed
field size $q$, and treats independent repeated observations. It does not
address correlated observations, hidden unknown source forms, or arbitrary
pairs of degenerate source and target forms.

\bibliographystyle{plain}
\bibliography{references}
\end{document}
